\originalchapter{宇宙热史、脱耦与热遗迹}
本章对应8.902第18讲末和第19讲。膨胀给一个时间尺度$H^{-1}$，反应给另一个时间尺度$\Gamma^{-1}$：比较它们，可以判断热平衡是否能跟上宇宙演化。本章先说明温度、熵和相对论自由度，再推导冷暗物质的冻结。每条方程都声明物理数密度还是共动量，以修正原课相关定义中的尺度因子问题。

\originalsection{黑体辐射的能量、数目与熵}
在频率$\nu$与$\nu+\dd\nu$内，单位体积含两个光子极化的态密度为$8\pi\nu^2\dd\nu/c^3$。每态平均占据数为$[\exp(h\nu/\kb T)-1]^{-1}$，故
\[
 u_\nu(T)=\frac{8\pi h\nu^3}{c^3}\frac1{\exp(h\nu/\kb T)-1},\quad
 u=\int u_\nu\dd\nu=\frac{\pi^2\kb^4}{15\hbar^3c^3}T^4.
\]
积分用$x=h\nu/(\kb T)$和$\int x^3/(\ee^x-1)\dd x=\pi^4/15$。各向同性光子动量通量给$p=u/3$。类似数目积分得
\[
 n_\gamma=\frac{2\zeta(3)}{\pi^2}\left(\frac{\kb T}{\hbar c}\right)^3,
 \qquad \frac{u}{n_\gamma}=\frac{\pi^4}{30\zeta(3)}\kb T\simeq2.701\kb T.
\]
光子的化学势在充分产生与吸收的热平衡中为零；对这种辐射，熵密度$s=(u+p)/T=4u/(3T)$。这里$s$有熵单位，后文取自然单位时可把$\kb$吸收入定义。

按频率与波长绘谱时最大点不同：$u_\nu$的峰满足$3(1-\ee^{-x})=x$，$x\simeq2.821$；每单位波长谱因雅可比$c/\lambda^2$变成$\lambda^{-5}/[\exp(hc/(\lambda\kb T))-1]$，峰满足$5(1-\ee^{-y})=y$，$y\simeq4.965$。不能用$\lambda=c/\nu_{\max}$把前者直接当成后者。

脱耦后无碰撞光子的相空间占据数沿轨道保持不变，物理动量$p\propto a^{-1}$。初始Planck分布$[\exp(pc/(\kb T_i))-1]^{-1}$仍可写成温度$T_i a_i/a$的Planck分布。因此\textbf{红移保持黑体形状并不要求每个时期仍在局部热平衡}。新的能量注入或不同温度混合可产生谱畸变；发生吸收与发射时也需另算碰撞项。

\originalsection{相对论自由度改变温度与时间关系}
本节起取$\hbar=c=\kb=1$，温度直接以能量表示，$1\ \mathrm{eV}$对应约$1.1605\times10^4$ K。热平衡物种的占据数为$f=(\exp[(E-\mu)/T]\mp1)^{-1}$，玻色取负、费米取正。化学势可忽略且$T\gg m$时
\[
 \epsilon=\frac{\pi^2}{30}g_*(T)T^4,\qquad
 s=\frac{2\pi^2}{45}g_{*s}(T)T^3.
\]
相同温度的每个费米自由度在能量与熵积分中相对玻色贡献$7/8$；不同温度物种在$g_*$中带$(T_i/T)^4$，在$g_{*s}$中带$(T_i/T)^3$。某组分变成非相对论性并湮灭后，它的熵转给仍耦合的粒子，不能维持所有物种的同一温度。

若没有净熵产生，$sa^3$守恒，故$T\propto g_{*s}^{-1/3}a^{-1}$。自由度不变时才简化为$T\propto a^{-1}$。辐射主导的Friedmann式给
\[
 H=1.66\sqrt{g_*}\frac{T^2}{M_P},\qquad M_P=G^{-1/2},
\]
这里$M_P$是\emph{非约化}Planck质量，与暴胀章$\bar M_P=M_P/\sqrt{8\pi}$区别明确。若$g_*$近恒定，$H=1/(2t)$，则
\[
 t=0.301\frac{M_P}{\sqrt{g_*}T^2}
 \simeq0.738\left(\frac{g_*}{10.75}\right)^{-1/2}
 \left(\frac{T}{\mathrm{MeV}}\right)^{-2}\mathrm s.
\]
自由度变化快或有非辐射主导阶段时不能用这一简单系数精确计时。

\originalsection{脱耦意味着反应不再能维持平衡}
每粒子的有效反应率$\Gamma$通常为$n\langle\sigma v\rangle$。当$\Gamma\gg H$且反应能改变相关变量时，系统可追踪平衡；当$\Gamma\ll H$，原有分布被膨胀带走而难以恢复。\emph{化学脱耦}冻结粒子总数，\emph{动理学脱耦}冻结与热浴的动量交换，两者可发生在不同时刻。

非相对论且非简并的平衡粒子有
\[
 n_{ \rm eq}=g\left(\frac{mT}{2\pi}\right)^{3/2}\exp[-(m-\mu)/T].
\]
当可与热浴湮灭且无不对称荷时取$\mu\simeq0$。降温后该密度随$\ee^{-m/T}$迅速降低，若湮灭一直足够快，粒子几乎消失；反应提前跟不上则留下遗迹。$\Gamma\simeq H$给数量级交叉点，真实偏离平衡时刻还取决于平衡丰度变化速率，不能将冻结当成瞬间关闭碰撞。

\originalsection{冷遗迹的Boltzmann方程与共动变量}
对稳定、有效自共轭的粒子，令$n$为物理数密度，反应截面约定使方程为
\begin{equation}\label{eq:astro-relic-boltzmann}
 \dot n+3Hn=-\langle\sigma v\rangle(n^2-n_{ \rm eq}^2).
\end{equation}
右侧的反向产生率由详细平衡确定。有粒子与反粒子分别计数、或并存多种共湮灭物种时，需明确相应有效截面和数目因子；本式不是任意粒子模型的全部动力学。

真正的共动粒子数密度是$N_c=na^3$。直接求导得
\[
 \dot N_c=-\frac{\langle\sigma v\rangle}{a^3}(N_c^2-N_{c,\rm eq}^2).
\]
原笔记的$n_c=n a^{-3}$及缺$a^{-3}$的共动碰撞式应修正。更便于热史的量是熵归一丰度$Y=n/s$；若$sa^3$守恒，膨胀稀释在$\dot Y$里自动消掉。

令$x=m/T$。在$g_{*s}$近恒定、辐射主导时$\dot x=Hx$，式\eqref{eq:astro-relic-boltzmann}变成
\begin{equation}\label{eq:astro-relic-y}
 \frac{\dd Y}{\dd x}=-\frac{s\langle\sigma v\rangle}{Hx}(Y^2-Y_{ \rm eq}^2),\qquad
 Y_{ \rm eq}\simeq0.145\frac g{g_{*s}}x^{3/2}\ee^{-x}.
\end{equation}
若熵自由度随温度变化，右侧还乘$1+Tg_{*s}'/(3g_{*s})$，因为$\dot x$相应减慢。额外熵注入或非标准膨胀则需重新建立两条背景方程。

\originalsection{冻结后的积分为何给出反比于截面的遗迹}
取$s$波湮灭、$\langle\sigma v\rangle$近恒定，定义
\[
 \lambda=0.264\frac{g_{*s}}{\sqrt{g_*}}M_Pm\langle\sigma v\rangle.
\]
则$s\langle\sigma v\rangle/(Hx)=\lambda/x^2$。冻结后$Y_{ \rm eq}$很小，式\eqref{eq:astro-relic-y}可积分为
\[
 \frac1{Y_\infty}=\frac1{Y(x_f)}+\frac\lambda{x_f}.
\]
若后续湮灭项主导，$Y_\infty\simeq x_f/\lambda$。所以更大的截面使粒子跟踪平衡更久、冻结温度更低，留下更少遗迹。只用$Y_\infty=Y_{ \rm eq}(x_f)$会漏掉冻结后的残余湮灭，通常不足以给可靠丰度。

常规弱相互作用量级的模型会给$x_f$约20--30，但值依赖质量、内部自由度和截面。今日密度为$\rho_{\chi,0}=m s_0Y_\infty$，对标准熵史换算有
\[
 \Omega_\chi h^2\simeq2.74\times10^8\left(\frac m{\mathrm{GeV}}\right)Y_\infty.
\]
数量级上，电弱尺度有效湮灭截面可产生所需宇宙学遗迹，这称WIMP机制的巧合；单个丰度不能唯一反推出$m$和$\sigma$，更不是已发现WIMP的证据。该例也不能排除非热产生、不对称暗物质或其他候选。

\originalsection{热、冷遗迹与自由流的尺度}
相对论性脱耦的粒子，在固定熵自由度下$Y$近恒定而不受Boltzmann指数抑制；非相对论性脱耦则受上述冻结控制。但热暗物质、暖暗物质、冷暗物质的结构效应最终取决于相关时期的\emph{速度分布和自由流尺度}，不能只看粒子第一次化学脱耦时是否相对论性。

无碰撞粒子的动量$p\propto a^{-1}$。速度$v=pc^2/\sqrt{m^2c^4+p^2c^2}$在相对论时期约为$c$，非相对论时期才随$a^{-1}$下降。共动自由流距离应积历史速度
\[
 \lambda_{\rm fs,c}(t)=\int_{t_{ \rm kd}}^t\frac{v(t')}{a(t')}\dd t'.
\]
原课用$v\sqrt{\pi/(G\rho)}$作Jeans类比给局部尺度，不能代替这个累计积分；其keV质量到几十Mpc的估计需指定热分布、组分温度和质量，不能推广到所有keV粒子。自由流擦除小尺度初始结构，大尺度可较早形成；冷暗物质则容许较小暗晕先出现并合并，但也存在微观截止，后期重子过程会改变可见星系。

\originalsection{热史中可靠计算与模型假设的分界}
核合成前后MeV温度的弱反应、中微子脱耦、电子正电子湮灭属于标准模型可计算过程。夸克到强子的转变在百MeV尺度，具体热力学不由理想气体式决定。更高温的大统一相变、Peccei--Quinn相变和Planck尺度量子引力阶段依赖相应理论是否正确，课程时间轴中的$10^{16}$、$10^{12}$ GeV不是已观测的宇宙事件。

光子热化还需区分弹性Compton散射与能改变光子数的双Compton、制动辐射等过程。原课“keV时光子脱离平衡并数目守恒”不是一个可由$T=m$判定的普遍瞬时事件。辐射—物质相等、复合及第一批天体形成也分别依赖背景密度、原子反应和结构演化。下一章利用这些明确的反应尺度推导轻元素核合成，而不把所有热史节点压成未经解释的时间表。
